% =====================================================================
% Paper 34 §III.28: Apparatus identity / state-side reduction
% =====================================================================
%
% Sprint 2, Track 9 (axis-completeness sub-sprint).  This subsection
% opens the state-side complement of Paper 34's projection dictionary.
% All twenty-seven prior §III entries (§III.1 Fock through §III.27 Wick
% rotation) operate on the spectral side --- the Dirac operator D, the
% Hilbert space H, the algebra A, the spectrum {lambda_n}.  None operate
% on the state side --- the density matrix rho on H, von Neumann entropy
% S(rho), reduced density matrices, mutual information, fidelity.
% §III.28 names that complement.
%
% Three-axis dictionary row (for §IV table integration):
%   Apparatus identity / state-side reduction
%     Variables added:  rho (density matrix) or beta = 1/T (Gibbs index)
%     Dimension:        preserves (S(rho) is dimensionless; T = energy/k_B
%                       when natural-unit conversion is consumed externally)
%     Transcendental:   von Neumann class (logarithm of operator);
%                       structurally distinct from master Mellin engine
%                       M1 / M2 / M3 (Sprint TD Track 5 PSLQ negative at
%                       100 dps against 12,312-form mechanical basis,
%                       coefficient ceiling up to 10^6)
%
% Cross-references in this subsection:
%   - Paper 27 (entropy-as-projection-artifact framing; S_full(GS) as
%     the T -> 0 anchor of the apparatus identity)
%   - §III.6 sec:proj_spectral_action (sibling on the OTHER side of the
%     spectral/state divide: spectral_action is the canonical D-side
%     Mellin-engine projection; apparatus_identity is its rho-side
%     complement)
%   - §III.24 sec:proj_adiabatic_BO (one of the rare projections that
%     creates the parameterized family rho_beta = e^{-beta H}/Z on which
%     state-side projections then act --- the slow/fast parameterization
%     of BO produces the family that the apparatus identity then samples)
%   - Paper 38 (Wasserstein-Kantorovich state space as the natural
%     domain for state-side metrics; sibling future-direction)
%
% Open-question flag to §VIII: the state-side complement of the
% dictionary is named but only this one entry currently populates it.
% Future Layer-2 state-side projections (mutual information,
% conditional entropy, fidelity, trace distance, Wasserstein-Kantorovich
% metric, relative entropy, quantum mutual information across bipartite
% splits) are candidate future entries.  This is a CANDIDATE STATE-SIDE
% ENUMERATION open question, distinct from existing spectral-side open
% questions on §VIII (calibration data, second packing axiom, etc.).
%
% =====================================================================

\subsection{Apparatus identity / state-side reduction}
\label{sec:proj_apparatus_identity}
\textbf{Source:} the framework's spectral-side output --- a Hamiltonian
$H$ on the Hilbert space $\mathcal{H}$ of a Fock-projected $S^3$ graph
(or its Bargmann--Segal / molecule-frame analogue), with eigenvalues
$\{\lambda_n\}$ and eigenstates $\{|\psi_n\rangle\}$ supplied by the
upstream projection chain.  All twenty-seven preceding entries of
\S~\ref{sec:layer2} (\S~\ref{sec:proj_fock} through the Wick-rotation
entry of Sprint~2) operate on this spectral side: they introduce
variables, dimensions, and transcendentals at the operator level of
$H$, $D$, or $\mathcal{H}$ itself.\\
\textbf{Target:} the state-side observables associated with a density
matrix $\rho$ on $\mathcal{H}$ --- specifically, the von Neumann entropy
$S(\rho) = -\mathrm{Tr}\,\rho \log\rho$, the reduced density matrices
$\rho_A = \mathrm{Tr}_B \rho_{AB}$ for any bipartite split
$\mathcal{H} = \mathcal{H}_A \otimes \mathcal{H}_B$, the mutual
information $I(A{:}B) = S(\rho_A) + S(\rho_B) - S(\rho_{AB})$, and the
classical--quantum apparatus identity
\begin{equation*}
   S_\mathrm{thermo}(T) \;=\; k_B \cdot S_\mathrm{microstate}(\rho_\beta),
   \qquad \rho_\beta = e^{-\beta H} / Z, \; \beta = 1/T,
\end{equation*}
which says that the macroscopic thermodynamic entropy of a system in
contact with a heat bath at temperature $T$ equals (in units of
$k_B$) the von Neumann entropy of the corresponding Gibbs density
matrix.  This projection is the rule that converts spectral-side
data into state-side observables.\\
\textbf{Variables introduced:} the density matrix $\rho$ on
$\mathcal{H}$, or, equivalently for thermal applications, the
inverse-temperature $\beta = 1/T$ that indexes the Gibbs family
$\rho_\beta = e^{-\beta H}/Z$.  The framework's spectral-side
machinery never references $\rho$; the apparatus identity is what
introduces it.  (When this projection follows
\S~\ref{sec:proj_observation}, $\beta$ is already present as the
Matsubara compactification radius; the apparatus identity then names
the state-side use of that $\beta$.  When it stands alone on a
non-thermal density matrix --- e.g.\ the ground-state reduced density
matrix $\rho_A = \mathrm{Tr}_B |\psi_0\rangle\langle\psi_0|$ of
Paper~27 --- $\rho$ is the new variable.)\\
\textbf{Dimension introduced:} preserves.  Von Neumann entropy
$S(\rho) = -\mathrm{Tr}\,\rho\log\rho$ is dimensionless by
construction (the eigenvalues of $\rho$ are pure probabilities and the
logarithm has no unit).  Temperature carries energy dimension via
$k_B T$, but the dimensional content of $T$ is consumed externally
through Boltzmann's constant rather than generated by the projection
itself; the projection's structural role is the dimensionless
operator-to-number map $\rho \mapsto S(\rho)$.\\
\textbf{Transcendental signature: von~Neumann class.}  The functional
$S(\rho) = -\mathrm{Tr}\,\rho\log\rho$ involves the logarithm of an
operator and is generically transcendental at the value level for any
non-trivial $\rho$.  Structurally, however, this transcendental class
is \emph{categorically distinct from the master Mellin engine}
$M_1 \cup M_2 \cup M_3$ that organises every spectral-side
transcendental documented in Paper~34 \S\S\ref{sec:proj_fock}--%
\ref{sec:proj_wick_rotation} and Paper~32 \S~VIII.  This is not a
framing claim --- it is a verified PSLQ negative.  Sprint~TD Track~5
(CLAUDE.md \S2, 2026-05-08) computed the framework's atomic
correlation entropies $S_\mathrm{full}(\mathrm{He}) =
0.040811051\ldots$, $S_\mathrm{full}(\mathrm{Li}^+) =
0.011211717\ldots$, and $S_\mathrm{full}(\mathrm{He},\, n_{\max}{=}4)
= 0.041879430\ldots$ to 150 dps, then ran PSLQ at 100 dps against a
$12{,}312$-form mechanical basis spanning the M1 (Hopf-base measure),
M2 (Seeley--DeWitt $\sqrt{\pi}$), and M3 (vertex-parity Hurwitz)
sub-rings plus algebraic and rational control classes.  Result: zero
hits across all three systems.  Stress tests at coefficient ceiling
$10^6$ against the ten simplest ring elements ($1, \pi, \pi^2, 1/\pi,
\zeta(2), \zeta(3), G, \log 2, \sqrt{2}, \varphi$) returned zero hits
at $10^{-50}$ tolerance.  Cross-system ratios
$S(\mathrm{He})/S(\mathrm{Li}^+) = 3.6400\ldots$ and the
$n_{\max}$-ratio $0.9744\ldots$ also null.  Six orders of magnitude in
coefficient ceiling do not change the verdict.  GeoVac correlation
entropy is genuinely a state-side von~Neumann quantity, structurally
disjoint from the spectral-side Mellin ring.\\
\textbf{Used in:} every place the framework computes an entropy or a
reduced density matrix.  Paper~27 (entropy as projection artifact;
$S_\mathrm{kin}/S_\mathrm{full} \sim 10^{-14}$ for He at
$n_{\max}=2,3$; $V_{ee}$-only generation of correlation entropy; the
universal $S_B = A(\tilde{w}_B/\delta_B)^\gamma$ scaling with
$\gamma_\infty \approx 1.96$ in the $Z, n_{\max} \to \infty$ limit)
exercises this projection throughout; the apparatus identity is the
structural mechanism that makes ``entropy'' a well-defined Layer-2
observable in the first place.  Sprint~TD Track~2 (apparatus identity
operational verification: $S_\mathrm{thermo} = k_B \cdot
S_\mathrm{microstate}$ to machine precision, $\max$ residual
$4.8\times 10^{-14}$, across H/He/Li$^+$ at $30$ log-spaced
temperatures each).  Future bipartite-entropy / mutual-information
analyses of atomic correlation (Paper~26 \S~VI Z-scaling
extrapolation).  Any future quantum-information-theoretic comparison
between the framework's qubit Hamiltonians (Paper~14) and competing
encodings, where fidelity and trace distance between candidate ground
states sit on the same state-side as $S(\rho)$.\\
\textbf{Empirical anchors.}  Two independent verifications, both at
machine precision.
\smallskip

\textit{Operational identity (Sprint~TD Track~2, 2026-05-08).}  The
classical-quantum apparatus identity $S_\mathrm{thermo}(T) = k_B \cdot
S_\mathrm{microstate}(\rho_\beta)$ is verified to floating-point
machine precision across the framework's hydrogenic-class systems.
Maximum residual across H/He/Li$^+$ at $30$ log-spaced temperatures
each: $4.8 \times 10^{-14}$.  The $T \to 0$ limit reproduces
Paper~27's ground-state correlation entropy $S_\mathrm{full}(\mathrm{GS})$
bit-identically: He $0.040811$ (residual $3.1\times 10^{-10}$),
Li$^+$ $0.011212$ (residual $3.2\times 10^{-10}$), H residual
identically zero by single-determinant ground-state rigidity.  The
$T \to \infty$ limit approaches $\log(\dim \mathcal{H})$ (e.g.\ $\log
14 \approx 2.639$ for hydrogen at $n_{\max}=3$), the Boltzmann maximum
mixing.  The Lindblad concavity bound holds at every $\beta$.  These
verifications confirm that the apparatus identity is operative as a
state-side reduction rule on the framework's hydrogenic systems
without anomaly.
\smallskip

\textit{Master-Mellin-engine disjointness (Sprint~TD Track~5,
2026-05-08).}  PSLQ at 100~dps against a $12{,}312$-form mechanical
basis (frozen before testing, W3-protocol style) returned zero hits
for $S_\mathrm{full}(\mathrm{He})$, $S_\mathrm{full}(\mathrm{Li}^+)$,
and $S_\mathrm{full}(\mathrm{He},\, n_{\max}{=}4)$.  Twenty random
comparable targets at coefficient ceiling $100$ also returned zero
hits (control null distribution genuinely zero).  Stress tests at
coefficient ceiling up to $10^6$ against the simplest ring elements,
and against cross-system ratios, also null.  The verdict is clean and
sharp: the framework's correlation entropies are not algebraic
combinations of the spectral-side master-Mellin transcendentals.\\
\textbf{Honest scope.}  State-side enumeration is currently
\emph{incomplete} in Paper~34's dictionary.  This entry is the first
state-side projection; the natural continuations
(mutual information, conditional entropy, fidelity, trace distance,
relative entropy $S(\rho \| \sigma)$, Wasserstein--Kantorovich
distance on the state space $P(\mathcal{H})$ of Paper~38's propinquity
machinery) are not enumerated here.  The structural status of those
candidate projections is named in the open-question section
(\S~\ref{sec:open}, candidate state-side enumeration); operationally,
they have appeared in the framework only as scaffolding inside
specific calculations (Paper~26 bipartite entanglement, Paper~38
Wasserstein--Kantorovich identification of the propinquity limit) and
have not yet been promoted to dictionary entries with their own
focal-length variable / dimension / transcendental tagging.  The
state-side dictionary should be read as deliberately opened in this
sprint, not as fully populated.

A second honest scope note: the apparatus identity is operationally
verified on the framework's hydrogen-class systems
(H, He, Li$^+$), where the spectral data is dense and the Gibbs
ensemble is small enough to enumerate.  Extension to molecular
spectral triples (Papers 15, 17, 19) and to the nuclear--electronic
composed triple (Paper~23 \S~VI) is mechanical but has not been
exercised; in those settings the relevant $\rho$ is potentially
high-dimensional and the von Neumann entropy may require sampling
or low-rank approximation methods rather than direct
diagonalisation.  This is an engineering scope rather than a
structural one: the projection is well-defined on every spectral
triple the framework produces; the operational cost of computing
$S(\rho)$ scales with the dimension of $\mathcal{H}$.

A third scope note concerns the relationship between this projection
and the foundational fact that pure-state von Neumann entropy of an
ungrouped system is zero.  Apparatus identity / state-side reduction
is non-trivial only when $\rho$ is mixed (Gibbs ensemble) or when a
bipartite trace $\rho_A = \mathrm{Tr}_B \rho_{AB}$ is taken on a pure
state.  This is consistent with Paper~27's framing: entropy is a
projection artifact in the precise sense that it requires choosing
either a thermal-bath partition (apparatus identity in the Gibbs
form) or a Hilbert-space partition (reduced density matrix).  The
state-side reduction is the structural operation that names that
choice and supplies the resulting dimensionless number.\\
\textbf{Structural reading.}  Categorically distinct from
\emph{every} prior entry of \S~\ref{sec:layer2}, in the precise
sense that those entries operate on the spectral data $(\mathcal{A},
\mathcal{H}, D)$ of the spectral triple, while this entry operates
on density matrices $\rho$ on $\mathcal{H}$.  The spectral / state
divide is the deepest taxonomic split in the dictionary so far: it
is sharper than the variable / dimension / transcendental axes
(Theorem~T3 in \S~\ref{sec:axiomatization}), sharper than
the master-Mellin sub-mechanism partition (M1 / M2 / M3), and sharper
than the Layer~1 / Layer~2 partition itself (which lives entirely on
the spectral side: Layer~1 is the bare graph $D$-data, Layer~2 is the
projected $D$-data).  The apparatus identity is structurally
\emph{the} sibling of \S~\ref{sec:proj_spectral_action} (Connes--%
Chamseddine spectral action): both project the framework's
finite-rank spectral data into a single dimensionless number, but
one does so via $\mathrm{Tr}\,f(D/\Lambda)$ on the spectral side
while the other does so via $\mathrm{Tr}\,\rho\log\rho$ on the
state side.  Sprint~TD Track~5's PSLQ negative is the quantitative
statement that these two procedures generate disjoint transcendental
classes: spectral-side traces of operator functions of $D$ live in
the master Mellin engine $M_1 \cup M_2 \cup M_3$ (Paper~32
\S~VIII case-exhaustion theorem; Sprint~TS-E3, May~2026), while
state-side traces of $\rho \log \rho$ do not.

The structural complement to \S~\ref{sec:proj_adiabatic_BO} is also
worth naming explicitly.  The Born--Oppenheimer projection creates a
parameterised family of Hamiltonians $H(\vec{R})$, and a thermal
ensemble of nuclear configurations weighted by Boltzmann factors
$e^{-\beta E(\vec{R})}$ is precisely the density matrix
$\rho_\beta(\vec{R}) = e^{-\beta H(\vec{R})}/Z$ on which the present
projection acts.  In the language of the multi-focal-composition arc
(May~2026, CLAUDE.md \S~2), adiabatic / BO supplies the Layer-1.5
parameter family; the apparatus identity is the Layer-2 reduction
that converts that family into a single thermodynamic number.  The
two projections compose naturally and the composition is exercised
operationally in Sprint~TD Track~2.

Paper~38's Wasserstein--Kantorovich state space $P(S^3)$
(Proposition~``Limit identification'') is the natural metric
container for the state-side complement of the dictionary.  In the
$n_{\max} \to \infty$ propinquity limit, the truncated spectral
triple $\mathcal{T}_{n_{\max}}$ converges to $\mathcal{T}_{S^3}$ with
the round-$S^3$ Wasserstein--Kantorovich distance as the metric on
its state space.  This identifies a clean future-direction target:
state-side projections beyond apparatus identity should sit naturally
in the Wasserstein--Kantorovich complex --- mutual information,
fidelity, and trace distance all have known metric structures on
state spaces of finite-rank operator algebras, and Paper~38's
machinery suggests these are reachable in the GH-convergent limit
without requiring additional NCG framework extension.  The state-side
dictionary is therefore named here as deliberately incomplete but
\emph{not} blocked: extension is structurally available and is the
natural next direction for the dictionary's axis-completeness.

A final structural reading: the apparatus identity / state-side
reduction is the projection that names ``how an observer extracts a
number from the system,'' as a structural complement to the
spectral-side projections that name ``how the framework constructs
the system in the first place.''  In the master-Mellin-engine
partition of Paper~32 \S~VIII and Paper~18 \S~III.7, the spectral
side has three sub-mechanisms (M1, M2, M3); the state-side --- as
documented by Sprint~TD Track~5's PSLQ negative --- is a fourth
mechanism class structurally outside the engine.  Whether the
state-side admits its own internal sub-mechanism partition
(analogous to the M1 / M2 / M3 split on the spectral side) is the
open structural question this entry leaves for future work.  Current
evidence (one verified entropy class, machine-precision apparatus
identity, PSLQ-disjoint from spectral side) is consistent with
either: a single state-side sub-mechanism that produces all
state-side observables uniformly, or multiple sub-mechanisms whose
distinctions are not yet operationally visible.  Either outcome would
sharpen the structural-skeleton-scope statement of CLAUDE.md \S~2
beyond its current spectral-side-only form.
